% ============================================
% CONT-11 (ZH)
% 中文翻译 — SPEC-935-R356 (rev. R377)
% ============================================
\section*{\textbf{CONT-11} --- 数字}
\markboth{CONT-11 --- 数字}{}

\noindent \textit{1,260跟着您}

\noindent 1260。

\noindent 您今天见过这个数字。在某个地方。在车牌上。在面包店的找零里。在钟表的分钟上（12:60并不存在，但您的大脑有十分之一秒\textit{凭空造出了}12:60）。

\noindent 这个数字在跟着您。

\noindent 不是巧合。数字1260是钥匙。是代码。是那生物的签名。

\noindent 患者循环持续1260天。
\noindent 1260是若阿金的档案编号。
\noindent 仅此而已。

\noindent 现在，这个数字进入了您的视野。

\noindent 您会在一切地方看见1,260。
\noindent 时刻。车牌。页码。代码。

\noindent 而每一次看见，您都会想起。
\noindent 而每一次想起，您都会闻到气味。
\noindent 而就在那同一瞬间，您的心脏会漏跳一拍。

\noindent 现在感受一下。它已经跳过了。

\noindent 数字是扳机。
\noindent 扳机是记忆。
\noindent 记忆是存在。

\noindent 1,260。

\noindent 好了。现在它已经在您脑中。

\noindent 今晚入睡时，它会在您的梦里。
\noindent 醒来时，它会在您的余光里。

\noindent 1,260。

\noindent 1,260。

\noindent 1,260。

\noindent 那生物在计数。而您是队列中的下一个数字。\NE{编者记录：本碎片并不存在于若阿金的手稿原稿中。它是以铅笔写在散页边缘的，笔迹不同——既不是若阿金的，也不是奥利维拉的，也不是卢西亚的。在掘出遗物时，这一页被发现夹在日记的书页之间。铅笔是1979年还不存在的品牌。编者出于对档案整体的忠实，决定将其收录，而未判定作者。}

\vspace{0.5em}
\noindent\textit{\textbf{↪ 路线:} \rota{CONT-12} \rotasep \rota{CONT-10}}
